% 出席確認クイズ 第10回　予測と分析の考え方
%   attendance/attendance10.html から生成（解説は本ガイド用に加筆）。

\quizq{1}{「相関があっても因果があるとは限らない」ことを示す例として適切なものはどれでしょうか?}%
  {因果関係があれば相関は必ずゼロになる}%
  {\correct{アイスの売上と水難事故が共に夏に増える(気温という別の要因がある)}}%
  {相関と因果はまったく同じ意味である}%
  {相関があれば必ず因果関係がある}%
  {正解は (b)。アイスの売上と水難事故はどちらも夏に増えるため相関がありますが、一方が他方の原因ではなく、気温という共通の要因（交絡因子）があるだけです。(a)(c)(d) は相関と因果の関係を誤解しています。経営データでも「広告費と売上」のように第三の要因（季節、景気）を疑う習慣が必要です。}

\quizq{2}{「教師あり学習」の説明として適切なものはどれでしょうか?}%
  {人がルールをすべて書く}%
  {ラベルのないデータをグループ分けする}%
  {\correct{正解ラベル付きのデータから入力と出力の関係を学習する}}%
  {報酬を最大化する行動を学ぶ}%
  {正解は (c)。教師あり学習は、正解ラベル付きのデータを使って、入力から出力を予測するモデルを学習する方法です（回帰・分類）。(b) は教師なし学習（クラスタリングなど）、(d) は強化学習、(a) はルールベースの説明です。}

\quizq{3}{「過学習(オーバーフィッティング)」の説明として適切なものはどれでしょうか?}%
  {データが多すぎて計算できない}%
  {学習が足りず精度が低い}%
  {\correct{学習データに合わせすぎて、新しいデータに対する精度が下がる}}%
  {モデルが小さすぎる}%
  {正解は (c)。過学習（オーバーフィッティング）は、学習データの細部や偶然のパターンまで覚え込んでしまい、未知のデータへの予測精度が落ちる現象です。(b) は逆に学習が足りない未学習（アンダーフィッティング）の説明です。}

\quizq{4}{予測モデルの精度を適切に評価する方法はどれでしょうか?}%
  {\correct{学習に使っていないデータ(テストデータ)で確かめる}}%
  {学習データでの精度だけを見る}%
  {精度は評価しない}%
  {モデルの大きさで判断する}%
  {正解は (a)。モデルの実力は、学習に使っていないテストデータでの精度で評価します。(b) 学習データでの精度だけを見ると過学習を見抜けません。交差検証（データを分割して評価を繰り返す方法）も広く使われます。}

\quizq{5}{AIの分析結果を経営で使う前に確認すべきこととして適切なものはどれでしょうか?}%
  {\correct{前提(使ったデータ・期間・仮定)と限界}}%
  {使ったパソコンの色}%
  {出力の文字数}%
  {AIが出したので確認は不要}%
  {正解は (a)。分析結果は、使用したデータ・期間・置かれた仮定によって変わるため、前提と限界を確認してから経営判断に使います。(d) AIが出した結果も検証が必要です。(b)(c) は分析の妥当性と無関係です。}
